\chapter{Tests et validation}
\label{ch:tests}

\section{Stratégie de test}

La stratégie retenue n'a pas visé une couverture maximale du code mais la protection des
règles métier dont une régression serait à la fois plausible et grave. Trois critères ont
guidé la sélection : la règle est-elle une exigence explicite du cahier des charges ?
Une erreur y serait-elle silencieuse ? Le code correspondant est-il susceptible d'évoluer ?

Cinq domaines satisfont ces trois critères et concentrent l'effort de test : la dérivation
du rôle effectif, le gel des comptes, l'arithmétique du stock de matériel, le ciblage des
notifications, et la matrice d'autorisation.

Deux niveaux de test coexistent. Les tests unitaires isolent un service au moyen de doubles
(Mockito) et s'exécutent sans contexte applicatif ni base de données ; ils sont rapides et
déterministes. Les tests d'intégration démarrent le contexte Spring complet sur une base
H2 en mémoire, avec le schéma dérivé des entités et les migrations désactivées ; ils
valident la chaîne de sécurité réelle. Ces derniers sont configurés pour ne jamais atteindre
la base de production : l'import du fichier de configuration local est explicitement
neutralisé dans le profil de test.

Le tableau~\ref{tab:suites} récapitule les suites.

\begin{table}[H]
\centering
\caption{Suites de tests automatisés}
\label{tab:suites}
\small
\begin{tabularx}{\textwidth}{@{}p{4.9cm}c p{2.2cm} X@{}}
\toprule
\textbf{Classe de test} & \textbf{Tests} & \textbf{Nature} & \textbf{Objet} \\
\midrule
\code{UserRoleSyncServiceTest} & 17 & Unitaire & Dérivation \srt{}~/~\rol{}~/~\man{}, réconciliation \\
\code{AuthorizationMatrixTest} & 9 & Intégration & Matrice d'autorisation des 9 rôles \\
\code{EquipmentServiceTest} & 8 & Unitaire & Invariant de stock du matériel \\
\code{NotificationServiceTest} & 5 & Unitaire & Ciblage des destinataires \\
\code{CustomUserDetailsTest} & 5 & Unitaire & Gel des comptes retraités et anciens \\
\code{BadRequestParamTest} & 4 & Intégration & Validation des paramètres de requête \\
\midrule
\textbf{Total} & \textbf{48} & & \\
\bottomrule
\end{tabularx}
\end{table}

\section{Tests unitaires de la dérivation des rôles}

C'est la suite la plus fournie, parce que la règle qu'elle protège est la plus structurante
du projet. Elle vérifie successivement :

\begin{itemize}
  \item la correspondance entre chacun des cinq statuts et son rôle dérivé, en l'absence de
        mandat ;
  \item la promotion par mandat, pour les trois types de mandats ;
  \item la priorité stricte entre mandats cumulés : un membre détenant simultanément un
        mandat de chef d'équipe et un mandat de directeur obtient le rôle de directeur ;
  \item la primauté du statut gelé : un membre retraité conservant par anomalie un mandat
        de directeur actif obtient malgré tout le rôle de retraité ;
  \item l'écriture systématique d'une ligne d'historique portant la cause du changement ;
  \item l'absence d'écriture lorsque le rôle calculé est déjà celui enregistré ;
  \item la correction d'une dérive par \code{reconcileAll()}, son \textbf{idempotence}
        --- une seconde exécution sur un système cohérent n'écrit rien --- et le
        complètement des historiques manquants.
\end{itemize}

Le test d'idempotence mérite d'être souligné : il garantit que la procédure de
réconciliation exécutée à chaque démarrage ne pollue pas l'historique en y ajoutant une
ligne à chaque redémarrage.

\section{Tests d'intégration}

Les deux suites d'intégration démarrent le contexte complet. Elles vérifient donc, au-delà
de la logique métier, que la configuration de sécurité, la chaîne de filtres, la
sérialisation JSON et le gestionnaire d'exceptions fonctionnent ensemble.

L'exécution de ces suites a d'ailleurs révélé un défaut qu'aucun test unitaire n'aurait
détecté, décrit en~§\ref{sec:defauts}.

\section{Tests d'autorisation}

\code{AuthorizationMatrixTest} vérifie la matrice du
tableau~\ref{tab:matrice} selon trois assertions systématiques par endpoint :

\begin{enumerate}
  \item une requête anonyme sur une ressource protégée reçoit un code~401 ;
  \item une requête portant un rôle insuffisant reçoit un code~403 ;
  \item une requête portant un rôle autorisé n'est jamais rejetée par 401 ni 403.
\end{enumerate}

La troisième assertion est volontairement formulée en négatif. Un rôle autorisé peut
légitimement recevoir un~404 --- l'entité visée n'existe pas dans la base de test --- ou un
400 ; ce qui importe est qu'il ne soit pas \emph{refusé}. Formuler l'assertion ainsi évite
de faire dépendre un test d'autorisation de la présence d'un jeu de données.

Les neuf tests couvrent les endpoints publics, une ressource ouverte à tout compte
authentifié, l'administration réservée à l'administrateur système, le journal d'audit et
les rapports réservés à la direction, les suppressions de niveau gouvernance, la fermeture
des modules internes aux doctorants, et l'accès explicite des doctorants à la création de
publications --- une autorisation qui pourrait sembler une erreur et qui est bien une
exigence du cahier des charges.

\section{Tests de validation des paramètres}

\code{BadRequestParamTest} verrouille une correction issue de l'audit : un paramètre de
requête absent ou inconvertible doit produire un code~400 et non un code~500. Les quatre
cas couverts sont un paramètre obligatoire manquant, une valeur d'énumération inconnue, un
identifiant UUID mal formé et un paramètre numérique non numérique.

\section{Tests du cycle de vie des membres}

Le gel des comptes est vérifié par \code{CustomUserDetailsTest} : un membre retraité ou
ancien n'est pas actif quelle que soit la valeur brute du champ correspondant en base ; un
compte technique dépourvu de profil de membre reste utilisable ; un compte
administrativement désactivé reste inactif indépendamment de son statut. La cohérence des
transitions elles-mêmes --- départ, retour, nouvelle période d'affiliation --- a été
validée par l'audit fonctionnel décrit en~§\ref{sec:e2e}.

\section{Tests des équipements}
\label{sec:testsequip}

\code{EquipmentServiceTest} vérifie l'invariant de stock sous huit angles, dont trois
méritent d'être cités : la distribution de l'intégralité du stock restant est autorisée et
amène la disponibilité à zéro ; une distribution excédant le stock est refusée \emph{sans
modifier l'état} ; un cycle complet de distribution puis de retour laisse le stock
exactement dans son état initial. Les deux derniers tests portent sur le rééquilibrage de
la disponibilité lors d'un changement de quantité totale, y compris le cas limite où la
diminution devrait produire une valeur négative --- le résultat est alors ramené à zéro.

\section{Tests des notifications}

\code{NotificationServiceTest} valide le ciblage. Le test le plus significatif vérifie que
la notification de la gouvernance atteint le directeur et le vice-directeur
\emph{actuellement en mandat}, destinataires résolus dynamiquement et non lus dans un champ
de rôle. Les autres cas couvrent la persistance d'une notification non lue, l'ignorance
silencieuse d'un destinataire inconnu, l'absence d'envoi lorsqu'aucun mandat n'est pourvu,
et la déduplication d'une liste de destinataires.

\section{Résultats des 48 tests}

L'exécution complète, sur le commit final, donne le résultat reproduit
au tableau~\ref{tab:resultats}.

\begin{table}[H]
\centering
\caption{Résultats d'exécution de la suite de tests}
\label{tab:resultats}
\small
\begin{tabularx}{\textwidth}{@{}p{5.2cm}p{2.8cm}X@{}}
\toprule
\textbf{Commande} & \textbf{Résultat} & \textbf{Détail} \\
\midrule
\code{mvn clean test} & \textbf{BUILD SUCCESS} & 48 tests exécutés, 0 échec, 0 erreur, 0 ignoré \\
\code{mvn clean compile} & \textbf{BUILD SUCCESS} & Compilation du backend sans avertissement bloquant \\
\code{npm run build} & \textbf{Succès} & 1\,919 modules transformés \\
\code{npx oxlint} & \textbf{0 erreur} & 6 avertissements préexistants sur l'export de constantes \\
\bottomrule
\end{tabularx}
\end{table}

Aucune mesure de couverture de code n'est rapportée : aucun outil de couverture n'a été
configuré sur le projet. Annoncer un pourcentage serait une invention. Ce que l'on peut
affirmer est plus limité et plus vérifiable : les cinq règles métier identifiées comme
critiques disposent chacune d'une suite dédiée, et la matrice d'autorisation est vérifiée
pour l'ensemble des rôles.

\section{Validation fonctionnelle de bout en bout}
\label{sec:e2e}

Les tests automatisés ont été précédés d'un audit fonctionnel manuel, conduit sur
l'application en fonctionnement. Il a porté sur les 21 routes de l'interface, les 9 rôles,
et 47 endpoints de l'API, chaque page étant parcourue et chaque commande déclenchée.

\subsection{Défauts détectés et corrigés}
\label{sec:defauts}

Cet audit a mis au jour quatre défauts réels, tous corrigés et, pour deux d'entre eux,
verrouillés par un test.

\paragraph{Dérive du rôle sur des comptes antérieurs.} Des comptes créés avant
l'introduction du service de dérivation conservaient un rôle obsolète, le service n'étant
invoqué que sur les chemins d'écriture. La correction a consisté à ajouter la procédure
\code{reconcileAll()} exécutée au démarrage, ainsi qu'une migration étendant la contrainte
de la colonne de cause pour admettre la valeur \code{RECONCILIATION}.

\paragraph{Commandes visibles mais refusées pour le vice-directeur.} La page de gouvernance
étant ouverte à trois rôles alors que la gestion des équipes et des mandats n'en autorise
que deux, un vice-directeur voyait des commandes que le serveur lui refusait. Les contrôles
concernés sont désormais masqués pour ce rôle, avec un encart expliquant ce qui lui reste
ouvert. Le test \code{deleteTeamIsAdminOrDirectorOnly} verrouille la règle côté serveur.

\paragraph{Paramètre de requête invalide traité en erreur serveur.} Un paramètre manquant
ou inconvertible aboutissait au gestionnaire d'exception générique et produisait un
code~500 au lieu d'un~400. Des gestionnaires dédiés ont été ajoutés ; la suite
\code{BadRequestParamTest} verrouille les quatre cas.

\paragraph{Corps de réponse malformé sur les erreurs d'authentification.} Défaut découvert
pendant l'écriture des tests d'autorisation, et invisible jusque-là. Les deux gestionnaires
d'erreur de la couche de filtrage sérialisaient leur réponse avec un convertisseur JSON
construit manuellement, dépourvu du module de gestion des dates. La sérialisation de
l'horodatage échouait : le client recevait bien un code~401, mais avec un corps vide, et
une exception était journalisée à chaque requête anonyme. La correction consiste à injecter
le convertisseur configuré par le contexte applicatif. Ce défaut illustre l'intérêt des
tests d'intégration : le code de statut, seul élément observé jusqu'alors, était correct.

\subsection{Résultats de la validation}

Le tableau~\ref{tab:e2e} synthétise les vérifications finales.

\begin{table}[H]
\centering
\caption{Résultats de la validation fonctionnelle de bout en bout}
\label{tab:e2e}
\small
\begin{tabularx}{\textwidth}{@{}p{6.2cm}X@{}}
\toprule
\textbf{Vérification} & \textbf{Résultat} \\
\midrule
Routes de l'interface parcourues & 21 / 21 rendues, aucun écran vide ni blocage \\
Gardes de route par rôle & Conformes : un doctorant est redirigé depuis les 8 routes qui lui sont fermées \\
Cohérence navigation / autorisation serveur & Conforme : la navigation d'un doctorant se limite à ses 6 modules \\
Dérive de rôle sur les comptes actifs & \textbf{0} --- le rôle dérivé coïncide avec le rôle stocké pour tous les membres \\
Réconciliation au démarrage & Idempotente : « tous les comptes sont déjà cohérents » \\
Comptes gelés & Authentification refusée (401) pour un statut retraité \\
Requête anonyme sur ressource protégée & 401 avec corps JSON complet \\
Exceptions non gérées côté serveur & 0 sur l'ensemble du parcours \\
\bottomrule
\end{tabularx}
\end{table}
